\subsection{File I/O}

\subsubsection{Read}

\noindent\textbf{Require}: Download the Cereal.csv file from the Canvas page and use the read.csv command to read in the csv file into R and assign it to the object called cereal.

\switchocg{W101}{\fbox{\textbf{Fold/Unfold}}}

\begin{ocg}{}{W101}{0}
\begin{lstlisting}[language=R]
library(readr)
library(dplyr)
library(ggplot2) 
library(here) # The `here` package simplifies the process of constructing file paths in R projects.

cereal <- read.csv(here("workshops/datasets", "Cereal.csv"))
cereal |> glimpse() # The `glimpse()` function from the `dplyr` package provides a quick overview of the `cereal` dataset.
\end{lstlisting}
\end{ocg}

\subsubsection{Data Frames}

\noindent\textbf{Require A}: Use head(cereal) to view the first few lines and class(cereal) to confirm it’s a data frame.

\switchocg{W102}{\fbox{\textbf{Fold/Unfold}}}

\begin{ocg}{}{W102}{0}
\begin{lstlisting}[language = R]
# using pipe 
cereal|>  head(7) 
\end{lstlisting}
\end{ocg}

\noindent\textbf{Require B}: Use colnames(cereal) for column names and dim(cereal) or nrow(cereal) for the number of rows.

\switchocg{W103}{\fbox{\textbf{Fold/Unfold}}}

\begin{ocg}{}{W103}{0}
\begin{lstlisting}[language = R]
cereal |> colnames
cereal |> nrow()
cereal |> dim()
\end{lstlisting}
\end{ocg}

\noindent\textbf{Require C}: Extract the calories column

\switchocg{W104}{\fbox{\textbf{Fold/Unfold}}}

\begin{ocg}{}{W104}{0}
\begin{lstlisting}[language = R]
# Extracts the 'calories' column from the 'cereal' data frame and assigns it to 'Cal' as a vector.
Cal <- cereal$calories 

# Selects the 'calories' column from the 'cereal' data frame and assigns it to 'Cal' as a data frame.
Cal <- cereal |> select(calories)

# Pulls the 'calories' column from the 'cereal' data frame and assigns it to 'Cal' as a vector.
Cal <- cereal |> pull(calories)
\end{lstlisting}
\end{ocg}

\noindent\textbf{Require D}: Extract rows 1 to 10 from the cereal data frame.

\switchocg{W105}{\fbox{\textbf{Fold/Unfold}}}

\begin{ocg}{}{W105}{0}
\begin{lstlisting}[language = R]
# Extracts the first 10 rows from the 'cereal' data frame using base R.
cereal[1:10,] 
\end{lstlisting}
\end{ocg}

\noindent\textbf{Require E}: Filter the cereal data frame to include only rows where mfr is "K" and assign it to Kelloggs.

\switchocg{W106}{\fbox{\textbf{Fold/Unfold}}}

\begin{ocg}{}{W106}{0}
\begin{lstlisting}[language = R]
Kelloggs <- cereal |> filter(mfr == "K") 
\end{lstlisting}
\end{ocg}

\subsubsection{Factors}

\noindent\textbf{Require A}: Load the Cereal data again with the read.csv command again. This time, use the optional argument, stringsAsFactors = TRUE. The mfr and type columns are now factors. Check that this is true.

\switchocg{W107}{\fbox{\textbf{Fold/Unfold}}}

\begin{ocg}{}{W107}{0}
\begin{lstlisting}[language = R]
# converting character columns to factors.
cereal.with.factors <- read.csv(here("workshops/datasets","Cereal.csv"), stringsAsFactors = TRUE)
cereal.with.factors
\end{lstlisting}
\end{ocg}

\noindent\textbf{Require B}: How many levels are there in mfr and type? (use the functions levels or nlevels) - levels will only count FACTORS not CHARACTER strings

\switchocg{W108}{\fbox{\textbf{Fold/Unfold}}}

\begin{ocg}{}{W108}{0}
\begin{lstlisting}[language = R]
levels(cereal.with.factors$mfr)
## or
nlevels(cereal.with.factors$mfr)
## or
str(cereal.with.factors$mfr)
\end{lstlisting}
\end{ocg}

\subsubsection{Vectors}

\noindent\textbf{Require A}: Extract the calories into a new vector called cereal.calories.

\switchocg{W107}{\fbox{\textbf{Fold/Unfold}}}

\begin{ocg}{}{W107}{0}
\begin{lstlisting}[language = R]
# Select the 'calories' column from the 'cereal' data frame and convert it to a vector.
cereal.calories <- cereal |> select(calories) |> pull()
\end{lstlisting}
\end{ocg}

\noindent\textbf{Require B}: How many elements are there in cereal.calories? (length)

\switchocg{W108}{\fbox{\textbf{Fold/Unfold}}}

\begin{ocg}{}{W108}{0}
\begin{lstlisting}[language = R]
# length(cereal.calories) - old way
cereal.calories |> length()
\end{lstlisting}
\end{ocg}

\noindent\textbf{Require C}: Extract the 5th to the 10th element from cereal.calories.

\switchocg{W109}{\fbox{\textbf{Fold/Unfold}}}

\begin{ocg}{}{W109}{0}
\begin{lstlisting}[language = R]
cereal.calories[5:10] 
\end{lstlisting}
\end{ocg}

\noindent\textbf{Require D}: Add one more element to cereal.calories using c().

\switchocg{W110}{\fbox{\textbf{Fold/Unfold}}}

\begin{ocg}{}{W110}{0}
\begin{lstlisting}[language = R]
# Add the value 1.0 to the end of the 'cereal.calories' vector using the concatenate function 'c'.
cereal.calories <- c(cereal.calories, 1.0) # c for concatenate

# Get the length of the 'cereal.calories' vector to see how many elements it contains.
length(cereal.calories)
\end{lstlisting}
\end{ocg}

\subsubsection{Matrix}

\noindent\textbf{Require A}: Can you force the cereal data frame to be a Matrix? (as.matrix(cereal)). Check that the elements have been forced into the character type.

\switchocg{W111}{\fbox{\textbf{Fold/Unfold}}}

\begin{ocg}{}{W111}{0}
\begin{lstlisting}[language = R]
# Convert the 'cereal' data frame to a matrix.
cereal.matrix <- as.matrix(cereal)

# Check if 'cereal.matrix' is indeed a matrix.
is.matrix(cereal.matrix)
\end{lstlisting}
\end{ocg}

\noindent\textbf{Require B}: Now do this again, but this time leave out the mfr, name and type columns. Check that the elements are now numeric.

\switchocg{W112}{\fbox{\textbf{Fold/Unfold}}}

\begin{ocg}{}{W112}{0}
\begin{lstlisting}[language = R]
# Remove the 'mfr', 'name', and 'type' columns from the 'cereal' data frame.
cereal.removed <- cereal |> select(-c(mfr, name, type))

# Display the column names of the modified 'cereal.removed' data frame.
cereal.removed |> colnames()

# Convert the 'cereal.removed' data frame to a numeric matrix.
cereal.numeric.matrix <- as.matrix(cereal.removed)

# Display the structure of the 'cereal.numeric.matrix'.
str(cereal.numeric.matrix)
\end{lstlisting}
\end{ocg}

\subsection{Numerical Summary}

\subsubsection{Summary}

\noindent\textbf{Require}:Use the summary function to extract the median, 1st quartile and 3rd quartile data from the sodium column.

\switchocg{W113}{\fbox{\textbf{Fold/Unfold}}}

\begin{ocg}{}{W113}{0}
\begin{lstlisting}[language = R]
summary(cereal$sodium)
\end{lstlisting}
\end{ocg}

\subsubsection{Basic Statistics}

\noindent\textbf{Require A}: Find the max, min, standard deviation and mean of the sodium (max(), min(), sd(), mean())

\switchocg{W114}{\fbox{\textbf{Fold/Unfold}}}

\begin{ocg}{}{W114}{0}
\begin{lstlisting}[language = R]
cereal |>  pull(sodium)|> summary()
\end{lstlisting}
\end{ocg}

\noindent\textbf{Require B}: Find the mean sodium of each mfr.

\switchocg{W115}{\fbox{\textbf{Fold/Unfold}}}

\begin{ocg}{}{W115}{0}
\begin{lstlisting}[language = R]
# Calculate the mean sodium content for each manufacturer ('mfr') in the 'cereal' data frame using the aggregate function.
mean.sodiums <- aggregate(sodium ~ mfr, data = cereal, FUN = mean)

# Display the resulting data frame with mean sodium values for each manufacturer.
mean.sodiums 

# Alternatively, select the 'sodium' and 'mfr' columns from the 'cereal' data frame, group by 'mfr', and calculate the mean sodium content for each group using the tidyverse approach.
cereal |>
    select(sodium, mfr) |>
    group_by(mfr) |>
    summarise(mean_sodium = mean(sodium))
\end{lstlisting}
\end{ocg}

\subsection{Graphical Summary}

\subsubsection{Boxplot}

\noindent\textbf{Require}: Make a boxplot of the sodium against mfr using boxplot().

\switchocg{W116}{\fbox{\textbf{Fold/Unfold}}}

\begin{ocg}{}{W116}{0}
\begin{lstlisting}[language = R]
# Create a boxplot of sodium content by manufacturer using base R.
# labels for both axes and give a meaningful title
boxplot(sodium ~ mfr, data = cereal,
        xlab = "Manufacturer", ylab = "Sodium content", main = "Sodium Content by Manufacturer")


# Using ggplot2.

cereal |> 
  ggplot(aes(x = mfr, y = sodium)) + 
  geom_boxplot() + 
  theme_classic() + 
  labs(x =  "Manufacturer", y = "Sodium content") +
  ggtitle("Sodium Content by Manufacturer")
\end{lstlisting}
\end{ocg}

\subsubsection{Scatter Plot}

\noindent\textbf{Require}: Plot calories against sodium using plot().

\switchocg{W117}{\fbox{\textbf{Fold/Unfold}}}

\begin{ocg}{}{W117}{0}
\begin{lstlisting}[language = R]
# Create a scatter plot of calories versus sodium using base R.
# The plot has the title "Calories versus Sodium Scatter Plot".
plot(calories ~ sodium, data = cereal, main = "Calories versus Sodium Scatter Plot")


# Create a scatter plot of calories versus sodium using ggplot2.
# Add a linear regression line without the confidence interval.
# The plot has the title "Calories versus Sodium Scatter Plot".
cereal |> 
  ggplot(aes(x = sodium, y = calories)) +
  geom_point() + 
  geom_smooth(method = "lm", se = FALSE) +
  labs(x =  "Sodium", y = "Carlories") +
  ggtitle("Calories versus Sodium Scatter Plot")
\end{lstlisting}
\end{ocg}

\subsection{Write Data to File}

\noindent\textbf{Require}: Write data frame with only the Kellogg’s observations to a file called kelloggs.csv. Use the write.csv command.

\switchocg{W118}{\fbox{\textbf{Fold/Unfold}}}

\begin{ocg}{}{W118}{0}
\begin{lstlisting}[language = R]
write.csv(Kelloggs, here("workshops/datasets", "kelloggs.csv"))
\end{lstlisting}
\end{ocg}